\section{REVISÃO DE LITERATURA E MARCO NORMATIVO}
\label{sec:revisao}

O objeto deste trabalho está entre duas literaturas que raramente se encontram. Uma delas trata do
que o direito exige da administração quanto a tornar inteligível --- e não apenas disponível --- o
modo como seus procedimentos correm; a outra, do instrumental capaz de reconstruir esse modo a
partir dos registros que os sistemas de tramitação já produzem. Nenhuma das duas, isoladamente, dá
conta do problema: a primeira pressupõe um relato do procedimento que não diz como produzir, e a
segunda produz esse relato sem indagar a que dever ele responde.

\subsection{Transparência, auditabilidade e procedimento administrativo}
\label{sec:rev-transparencia}

A literatura sobre transparência administrativa abandonou em larga medida a equação entre
divulgação e \emph{accountability}. \citeonline{ananny2018seeing} expõem por que ver um sistema não
equivale a conhecê-lo nem a governá-lo, e \citeonline{janssen2012benefits} encontram a mesma
confusão no governo aberto, em que publicação é rotineiramente tomada por uso.
\citeonline{bovens2007analysing} fornece o vocabulário de que a distinção precisa:
\emph{accountability} é uma relação em que um ator explica sua conduta a um foro competente para
questionar e julgar. A doutrina administrativista brasileira formula o mesmo ponto em termos mais
exigentes: transparência é ``publicidade qualificada'', e ``não basta que a informação esteja
disponível se ela não for inteligível'' \cite[p.~245]{brbaptista2023transparencia}; à publicação
meramente formal ``falta a visibilidade real e a inteligibilidade geral da ação administrativa''
\cite[p.~107--108]{brrodrigues2014publicidade}, que distingue a publicidade proativa --- a
transparência ativa --- da reativa, acionada por pedido. Um catálogo não é, sob essa leitura, uma
divulgação, mas o conteúdo explicativo que tal foro requer. A mesma doutrina chega a enunciar a tese
normativa deste trabalho sem qualquer referência a mineração de processos: ``não se pode conceber a
publicidade sem que se materialize o íter procedimental percorrido pelos agentes... até a edição do
ato administrativo'' \cite[p.~255--256]{brbaptista2023transparencia}.

Há literatura que trata o procedimento administrativo nesse mesmo nível.
\citeonline{suksi2021administrative} argumenta que as salvaguardas de devido processo escritas para
servidores humanos se diluem à medida que o procedimento é automatizado, e que as que valem a pena
são preventivas --- argumento para tornar inspecionável o procedimento tal como ele é antes de
automatizar qualquer coisa sobre ele. \citeonline{branting2025sensitive} tratam registros
governamentais sujeitos a obrigações de acesso à informação, e \citeonline{tinarrage2025empirical}
avaliam empiricamente um instituto do direito brasileiro associando resultados de classificação a
uma leitura jurídica de seu desempenho. Próximo em método, \citeonline{vianna2024topic} descobrem
estrutura temática sem supervisão sobre documentos jurídicos em português brasileiro e submetem o
resultado a especialistas em direito --- etapa pela qual o catálogo aqui produzido não passou (seção
4.11). Duas restrições decorrem daí: o que os registros contêm decide qual tarefa pode ser tentada
sobre eles \cite{medvedeva2023rethinking}, e uma taxonomia oferecida para uso jurídico é mais
defensável quando produzida por procedimento declarado, com condições explícitas de encerramento
\cite{nickerson2013taxonomy}.

\subsection{Marco normativo brasileiro}
\label{sec:rev-normativo}

\subsubsection{Princípios constitucionais e garantias do art. 5\textordmasculine{}}

A administração pública brasileira está vinculada, em todos os níveis da federação, aos cinco
princípios do \emph{caput} do art.~37 da Constituição de 1988 \cite{brasil1988cf}: legalidade,
impessoalidade, moralidade, publicidade e eficiência, este último inserido pela Emenda
Constitucional n\textordmasculine{}~19/1998. São princípios, e não regras, e a distinção decide
para que um registro de atividades pode ser usado. A leitura dos princípios como mandados de otimização é
corrente na retórica decisória brasileira, mas é expressamente contestada como descrição do que a
doutrina administrativista nacional de fato faz \cite{brgabardo2017eficiencia}, e nada do que segue
depende de como a controvérsia se resolva: apenas a consequência probatória interessa aqui. Uma
regra é cumprida ou descumprida; um princípio admite mais e menos. As grandezas contínuas que um registro de atividades
produz --- tempo decorrido, intervalos ociosos, contagens de repasse, repetição --- correspondem a
um gradiente, não à qualificação jurídica caso a caso que uma constatação de ilicitude exigiria.
Avaliar instrumentos de governo contra princípios já vigentes tem precedente
\cite{oswald2018algorithm}, e nenhuma afirmação deste trabalho se apoia em dispositivo de prazo.

O princípio da eficiência pesa aqui mais do que um lema constitucional pesaria, e a doutrina
brasileira diz por quê em termos que incidem sobre um registro de atividades. A leitura consagrada distingue a
eficiência no modo de atuação do agente público daquela no \emph{modo de organizar, estruturar e
disciplinar} a Administração, e é este segundo aspecto que governa o procedimento, ``dada a
necessidade de previsibilidade da atuação administrativa''
\cite[p.~180]{brchicoski2004eficiencia}. Dele a mesma literatura extrai deveres reconhecíveis em um
registro de atividades: a celeridade veda ``paralisações desnecessárias'' e o ``ritmo lento''; a economia processual
veda a ``repetição desnecessária de atos ou fases''; a oficialidade impede que a Administração
aguarde reiterados pleitos das partes para praticar atos de sua competência
\cite[p.~184--185]{brchicoski2004eficiencia}. A eficiência é tratada como norma constitucional de
eficácia direta, e entre os controles que autoriza incluem-se os métodos de trabalho e o próprio
procedimento \cite{brgabardo2017eficiencia,brfranca2000eficiencia}. Intervalos ociosos, laços de
retrabalho e processos parados são, portanto, a forma observável de questão que essa literatura já
trata como jurídica.

Duas garantias do art.~5\textordmasculine{} acrescentam o lado voltado ao cidadão: o inciso LXXVIII,
da Emenda Constitucional n\textordmasculine{}~45/2004, assegura a todos, no processo judicial e no
administrativo, a razoável duração e os meios que garantam a celeridade da tramitação; o inciso
XXXIV, alínea ``a'', garante, independentemente de taxas, o direito de petição em defesa de direitos
ou contra ilegalidade ou abuso de poder. Um direito de petição vale o que o peticionário consegue
descobrir, e saber que um procedimento existe não é saber como ele corre, quanto tempo leva ou por
quais unidades passa.

\subsubsection{A lei estadual de processo administrativo e a transparência}

No nível estadual, a lei que regula o processo administrativo no âmbito da Administração Pública
do Estado do Piauí --- e, portanto, no âmbito da SEAD-PI --- enuncia um rol de princípios
materialmente mais amplo que o de sua congênere federal, e o abre pela transparência:

\begin{citacao}
O processo administrativo obedecerá, dentre outros, aos princípios da transparência, legalidade,
finalidade, motivação, razoabilidade, proporcionalidade, moralidade, ampla defesa, contraditório,
segurança jurídica, impessoalidade, eficiência, celeridade, oficialidade, publicidade, participação,
proteção da confiança legítima e interesse público.
\cite[art.~2\textordmasculine{}]{state2016procedurelaw}
\end{citacao}

\noindent A posição da transparência à frente do rol não é decorativa, e a oficialidade carrega
peso próprio: onde a Administração deve impulsionar o procedimento por iniciativa própria, o tempo
parado é fato administrativo, e não demora imputável à parte.

\subsubsection{O dever de publicar as principais etapas: LAI e Lei 13.460/2017}

Duas leis convertem a publicidade em dever operacional, e cada uma fornece coisa diversa. A
Lei de Acesso à Informação \cite{brasil2011lai} faz da publicidade a regra e do sigilo a exceção, e
não condiciona a divulgação a requerimento: seu art.~8\textordmasculine{} impõe a divulgação
espontânea de informações de interesse coletivo ou geral, e o \S~3\textordmasculine{}, III exige dos
sítios oficiais o acesso automatizado por sistemas externos em formatos abertos, estruturados e
legíveis por máquina. Esse dispositivo governa o \emph{canal e o formato}, e não o que deve ser
descrito --- o rol de conteúdo mínimo do \S~1\textordmasculine{} não contém item algum sobre etapas
de tramitação.

O dispositivo que o diz é o art.~7\textordmasculine{}, \S~2\textordmasculine{}, III da lei que rege
os usuários dos serviços públicos:

\begin{citacao}
Art.~7\textordmasculine{} Os órgãos e entidades abrangidos por esta Lei divulgarão Carta de Serviços
ao Usuário. {[...]} \S~2\textordmasculine{} A Carta de Serviços ao Usuário deverá trazer informações
claras e precisas em relação a cada um dos serviços prestados, apresentando, no mínimo, informações
relacionadas a: {[...]} III --- principais etapas para processamento do serviço; {[...]}
\S~4\textordmasculine{} A Carta de Serviços ao Usuário será objeto de atualização periódica e de
permanente divulgação mediante publicação em sítio eletrônico do órgão ou entidade na internet.
\cite[art.~7\textordmasculine{}]{brasil2017cdu}
\end{citacao}

\noindent O dever é proativo, o canal eletrônico é obrigatório, a atualização é periódica e a
enumeração é expressamente um piso --- ``no mínimo'' ---, não um teto. O órgão federal que orienta
a transparência ativa arrola a Carta de Serviços entre os conteúdos mínimos obrigatórios dos sítios
institucionais, remetendo seu conteúdo ao regulamento federal da matéria
\cite{brcgu2019gta,brbrasil2017decreto9094}: a ponte entre o dever de publicar etapas e o regime da
Lei de Acesso à Informação é administrativa, e não doutrinária. Lidos em conjunto, os dois diplomas
equivalem a um dever de publicar um relato de como um procedimento efetivamente corre, em forma que
uma máquina possa consumir. É desse dispositivo que o trabalho depende, e a razão é uma
pressuposição, não uma proibição: publicar as principais etapas pressupõe saber quais são essas
etapas, e nenhuma das duas leis diz como esse conhecimento deve ser produzido nem supõe que ele
possa faltar.

\subsubsection{Anonimização e o art. 12 da LGPD}

A proteção de dados constringe o mesmo material pelo lado oposto. O art.~12 da Lei Geral de
Proteção de Dados Pessoais coloca o dado anonimizado fora do alcance da lei, com ressalva de
reversibilidade:

\begin{citacao}
Os dados anonimizados não serão considerados dados pessoais para os fins desta Lei, salvo quando o
processo de anonimização ao qual foram submetidos for revertido, utilizando exclusivamente meios
próprios, ou quando, com esforços razoáveis, puder ser revertido.
\cite[art.~12]{brasil2018lgpd}
\end{citacao}

\noindent Os arts.~5\textordmasculine{}, III e 5\textordmasculine{}, XI definem a anonimização por
referência a meios técnicos razoáveis e disponíveis na ocasião do tratamento. É padrão de risco, e
não estado binário, o que converte a robustez de uma anonimização em questão empírica sobre a
técnica disponível --- razão pela qual este estudo a trata como braço experimental cujo custo
a jusante é medido, e não como detalhe de pré-processamento. O \S~2\textordmasculine{} opera em
sentido inverso e delimita um risco: dados utilizados para formação do perfil comportamental de
pessoa natural identificada poderão igualmente ser considerados pessoais, cláusula capaz de
reconduzir traços por caso ao âmbito da lei a despeito do \emph{caput}.

\subsubsection{O alcance do dever sobre este corpus: uma questão em aberto}
\label{sec:rev-alcance}

O alcance do dever de publicar as principais etapas sobre o corpus estudado é o limite mais agudo da
afirmação normativa deste trabalho, e é questão em aberto. A Lei n\textordmasculine{}~13.460/2017
vincula seus deveres a
serviço público prestado a \emph{usuário} --- quem se beneficia ou utiliza, efetiva ou
potencialmente, de serviço público, na definição de seu art.~2\textordmasculine{}, I ---, e o corpus
aqui analisado é predominantemente procedimento administrativo interno cuja parte é um servidor
público (seção 3.3). Se um servidor que transaciona com sua própria administração é usuário nesses
termos é questão de direito administrativo brasileiro que este trabalho não resolve, sobre a qual
nada aqui deve ser lido como tomada de posição, e de cuja resposta nenhuma afirmação deste trabalho
depende.

Há leitura séria de que é. A definição do inciso I é funcional, e não estatutária: fala de quem
\emph{se beneficia ou utiliza} e não exclui quem seja também agente da Administração --- o servidor
que requer licença ou contesta um desconto em folha está na ponta receptora de um serviço, e não no
exercício de autoridade, com a assimetria que a lei existe para corrigir. O inciso II a reforça ao
definir serviço público como ``atividade administrativa \emph{ou} de prestação direta ou indireta de
bens ou serviços à população'', pois a disjunção admite
que ``atividade administrativa'' seja limbo autônomo, não qualificado por ``à população'' --- no que
a acompanham o manual do órgão federal de ouvidoria, para o qual ela é o ``ato voltado para
organização e funcionamento das estruturas estatais'' \cite[p.~68]{brcgu2018ouvidoria}, e a doutrina
que ali vê os serviços públicos administrativos, inclusive os prestados por um órgão público a outro
\cite{brmarrara2018cdusp}. A norma federal que estende aos agentes públicos internos ``as mesmas
garantias e proteções destinadas aos usuários'' \cite[art.~48]{brcgu2024portaria116} milita no mesmo
sentido, pois um regulador não tem competência para estender as proteções de uma lei a uma classe
que a lei não alcança de modo algum.

Há leitura séria de que não é. ``À população'' pode governar toda a enumeração do inciso II, que
então se lê em conjunto com o inciso I e com o objeto declarado no art.~1\textordmasculine{},
\emph{caput} --- o usuário de um serviço ---, e o mesmo autor que sustenta a leitura ampla delimita
o Código à função administrativa prestativa \cite{brmarrara2018cdusp}, o que a estreita justamente
onde a Administração nada presta ao servidor. O art.~2\textordmasculine{}, IV opõe objeção textual
ao definir \emph{agente público} como categoria distinta do usuário do inciso I, e a resposta
disponível --- que os dois incisos classificam transações, e não pessoas, de modo que a mesma pessoa
ocupa as duas posições em momentos diferentes --- é argumento, e não doutrina assentada. A mesma
norma de ouvidoria corta também aqui, pois uma norma que precisa \emph{equiparar} agentes internos a
usuários pressupõe que eles já não o sejam, e a equiparação foi editada para o regime de ouvidoria,
não para o dever de publicação. E o órgão federal tampouco é uniforme: seu guia de transparência
ativa restringe o dever de publicar a Carta aos órgãos ``que prestam serviços diretamente ao
público'' \cite{brcgu2019gta}. Lida amplamente, a lei alcançaria parcela substancial deste corpus;
lida estritamente, apenas a minoria voltada para fora.

Nenhum instrumento fecha a questão. O decreto que regulamenta a Carta de Serviços no Poder Executivo
federal exige que dela constem ``as etapas, presentes e futuras, esperadas para a realização dos
serviços, incluídas a estimativas de prazos'' e ``os mecanismos para a consulta pelos usuários
acerca das etapas, cumpridas e pendentes''
\cite[art.~11, \S~3\textordmasculine{}]{brbrasil2017decreto9094} --- descrição do artefato produzido
neste trabalho mais próxima do que qualquer coisa na própria lei ---, mas seu
art.~1\textordmasculine{} vincula \emph{os órgãos e as entidades do Poder Executivo federal}, e o
ato sobre ouvidorias tem o mesmo recorte: diante de uma secretaria estadual, ambos são persuasivos,
não vinculantes. Nenhuma decisão dos tribunais superiores ou do órgão federal de controle externo
parece resolvê-la, e há explicação documentada para o silêncio: levantamento em tribunal estadual
entre 2018 e 2020 encontrou apenas duas decisões que sequer mencionaram a lei, e o Código de Defesa
do Consumidor foi aplicado em catorze de quinze acórdãos sobre prestação inadequada de serviço
público \cite{brjuruena2021usuario}. A questão é viva o bastante para o regulador federal tê-la
enfrentado, e permanece aberta.

Dois outros fundamentos incidem sobre o artefato sem depender dela, e ambos são mais fracos que o
gancho da Carta de Serviços: o dever de publicidade proativa da Lei de Acesso à Informação
\cite{brasil2011lai} e o princípio da transparência que abre o art.~2\textordmasculine{} da lei
estadual de processo administrativo \cite{state2016procedurelaw}, que vincula a condução dos
procedimentos como tais e não se dirige a classe alguma de beneficiário. Nenhum deles gera dever de
\emph{produzir} o artefato: o primeiro se prende a informação que o órgão produz ou custodia, e o
registro cuja ausência é a premissa deste trabalho não é, por hipótese, nem uma coisa nem outra, de
modo que o dever só o alcança depois que ele existe; e um princípio satisfeito em graus não pode
compelir instrumento determinado com mais facilidade do que pode render a constatação de ilicitude
caso a caso que a seção 2.2.1 recusou. O que sustentam é mais estreito: existindo o catálogo,
publicá-lo deixa de ser discricionário como seria publicar um documento interno de trabalho, e há
razão, admitindo graus, para preferir uma administração capaz de prestar contas de como seus
procedimentos correm. Produzi-lo é, portanto, prudencial e normativo, e não compelido, exceto onde o
dever da Carta de Serviços incida e forneça a especificação faltante do que deve ser descrito.

O análogo conceitual mais próximo fora do Brasil --- o direito à boa administração do art.~41 da
Carta dos Direitos Fundamentais da União Europeia, que exige tratamento imparcial, equitativo e em
prazo razoável --- não se aplica aqui: dirige-se, por seus próprios termos, às ``instituições,
órgãos e organismos da União'' \cite{eu2012charter}, e o Tribunal de Justiça decidiu que não alcança
os Estados-Membros \cite[par.~67]{cjeu2014ys}, leitura depois confirmada
\cite[par.~44]{cjeu2014mukarubega}, ainda que os mesmos julgados preservem a boa administração como
princípio geral do direito da União. A menção localiza uma ideia compartilhada, e não uma fonte de
obrigação aplicável à administração estudada.

Lidos em conjunto, esses dispositivos pressupõem algo que não fornecem: um registro de como a
administração efetivamente opera. \citeonline{brauneis2018algorithmic} testaram essa pressuposição
empiricamente, protocolando 42 pedidos de acesso a registros em 23 estados norte-americanos sobre
programas governamentais de algoritmos preditivos, e constataram que o obstáculo recorrente não era
sigilo comercial, mas o fato de que os registros nunca haviam sido criados. A premissa deste
trabalho é carregada por esse único estudo de outra ordem jurídica, e sua força é limitada: aqueles
pedidos versavam sobre matéria nova e mediada por contratação, em regime acionado por requerimento e
sob leis com exceções de segredo comercial, ao passo que o cenário presente difere nos quatro pontos
--- procedimento de rotina, dever de publicação proativa, ausência de camada de fornecedor e
princípio constitucional de publicidade. A inferência de que a ausência de registro se repete aqui é
extrapolação não testada: não se detém evidência de que a administração estudada carecesse de um
relato documentado de seus procedimentos antes deste trabalho, e estabelecê-la exigiria pedido de
acesso, consulta direta ao órgão ou a literatura empírica sobre cumprimento da legislação de acesso
à informação, nenhum deles empreendido aqui. A evidência brasileira mais próxima é de outra
natureza: auditoria de portais estaduais e municipais encontrou buscas por Carta de Serviços que
frequentemente não retornaram resultado algum \cite{brjuruena2021usuario}, descumprimento do dever
de publicar que não é o mesmo que registro nunca produzido.

A contribuição é enunciada nesse nível, e não acima dele. Um catálogo dos tipos de processo tal como
eles são, cada um com seu caminho típico derivado dos registros de eventos da própria administração,
fornece um desses registros: enunciado das principais etapas no sentido do
art.~7\textordmasculine{}, \S~2\textordmasculine{}, III, legível por máquina no sentido do
art.~8\textordmasculine{}, \S~3\textordmasculine{}, III, regerável a partir do registro de atividades
e inspecionável pelos servidores que operam os procedimentos. Os tipos são descritivos e revisáveis,
não classificações normativas, e o catálogo não é prova de ilegalidade, constatação de conformidade
nem base para medida contra servidores identificados.

\subsubsection{Do dever de transparência ao dever de otimizar}
\label{sec:rev-otimizar}

O enquadramento até aqui tratou o catálogo como artefato de transparência. Há, porém, um segundo
fundamento, e ele é de lei ordinária. A Lei n\textordmasculine{} 14.129, de 2021, enuncia como
próprio objeto ``o aumento da eficiência da administração pública'' e arrola, entre os princípios
e diretrizes do Governo Digital, a transparência na execução dos serviços públicos somada ao
monitoramento de sua qualidade e --- na formulação mais próxima do que aqui se produz --- ``o uso
da tecnologia para otimizar processos de trabalho da administração pública''
\cite{brasil2021govdigital}. Trata-se de enunciado legislativo de que se espera da administração
o emprego de tecnologia sobre os seus próprios procedimentos de trabalho, e não apenas sobre o
serviço que presta para fora. Seu alcance é delimitado como o das demais normas federais
discutidas acima: atinge os demais entes federados apenas na medida em que estes adotem seus
comandos por atos normativos próprios, de modo que, perante a administração estudada, é
persuasivo e não vinculante. A norma que vincula é a estadual, cujo art.~2\textordmasculine{}
abre o rol de princípios com a \emph{transparência} e nele inclui \emph{eficiência},
\emph{celeridade} e \emph{oficialidade} \cite{state2016procedurelaw}.

Disso decorre uma consequência distinta do argumento de transparência. Se otimizar
procedimentos é dever, e não preferência gerencial, então descrever como esses procedimentos de
fato correm não é somente ato de divulgação: é condição de possibilidade do próprio dever, pois
não se aperfeiçoa aquilo que não se conhece. Os caminhos típicos aqui produzidos são precisamente
o objeto sobre o qual tal análise incidiria --- se o caminho observado é o que o procedimento
deveria seguir, onde ele se desvia, que etapas poderiam ser suprimidas --- e essa análise é
trabalho futuro que este estudo viabiliza, e não realiza.

\subsubsection{Por que o registro não pode ser presumido: o sistema não restringe o percurso}
\label{sec:rev-naorestringe}

Um último ponto é técnico antes de ser normativo, e sustenta a escolha de \emph{derivar} o
catálogo em vez de transcrevê-lo de uma classificação existente. O SEI registra ações e
tramitações, mas não impõe, por tipo de processo, qual sequência de atividades é admissível. Nada
no sistema impede que dois processos administrativos do mesmo tipo declarado percorram trajetos
inteiramente distintos, nem obriga que um tipo declarado corresponda a um procedimento
determinado. A variação, portanto, não é anomalia a ser explicada: é o comportamento esperado de
um sistema que não constrange o percurso.

A consequência é estrutural. Onde o procedimento não é restringido pelo sistema nem descrito em
outro lugar, o tipo declarado não serve como prova do que foi feito, e o único relato do
\emph{as-is} em que se pode confiar é o que se deriva do registro daquilo que efetivamente
ocorreu. É essa lacuna que o presente trabalho endereça, e é ela que torna a taxonomia declarada
uma hipótese a ser testada --- como se faz na seção 4.1 --- e não um dado de partida.

\subsection{Mineração de processos e clusterização de traços}
\label{sec:rev-clustering}

Minerar um único modelo a partir de um registro de atividades heterogêneo produz uma estrutura que não descreve
comportamento algum efetivamente observado, e o remédio consolidado é particionar o registro de atividades primeiro e
descobrir um modelo especializado para cada grupo \cite{greco2006clustering}. Particionar exige uma
representação: perfis de traço sobre atividades, transições, executores e atributos do caso tornam
aplicável o agrupamento em espaço vetorial, e perfis distintos produzem partições materialmente
distintas \cite{song2009trace}; subsequências conservadas restauram a informação de ordem que os
perfis de sacola de atividades descartam \cite{bose2009context}; representações densas aprendidas
superam perfis manuais \cite{dekoninck2018act2vec}. Uma revisão sistemática organiza quinze anos
desse trabalho em quatro dimensões --- representação, algoritmo, uso de informação fora do fluxo de
controle e estratégia de avaliação --- e já arrola a qualidade do modelo a jusante entre as
estratégias de avaliação reconhecidas \cite{zandkarimi2020generic}. A instância mais nítida é o
ActiTraC, que só admite um traço em um grupo quando a aderência do modelo descoberto para aquele
grupo não se degrada significativamente: é o artefato a jusante que conduz o agrupamento, e não uma
distância no espaço de atributos \cite{deweerdt2013actitrac}. O que essa literatura chama de
variante aproxima-se do que aqui se chama tipo de processo \cite{taymouri2021variant}.

A diferença deste trabalho em relação a essa linha é de posicionamento do escore a jusante, e não de
princípio; a seção~\ref{sec:desenho} a especifica. O precedente mais próximo em técnica e jurisdição ao mesmo
tempo é o de
\citeonline{araujo2025trace}, que aplicam clusterização de traços a dados de processos judiciais
brasileiros para simplificar esses procedimentos e identificar padrões neles. A diferença de cenário é
substantiva: o procedimento judicial é codificado externamente, e uma partição dele pode ser
conferida contra um código de processo, ao passo que a tramitação administrativa aqui estudada não
tem especificação externa correspondente --- razão pela qual o registro precisa ser derivado, e não
consultado.

\subsection{Modelos sequenciais por grupo, caminhos representativos e monitoramento preditivo}
\label{sec:rev-sequencia}

Treinar um modelo sequencial por grupo não é arranjo recente. \citeonline{ferreira2007sequence}
tratam o agrupamento como mistura de cadeias de Markov de primeira ordem, de modo que cada grupo
\emph{é} um modelo sequencial e um traço pertence ao grupo cujo modelo mais provavelmente o gerou;
\citeonline{difrancescomarino2019clusteringppm} separam as duas etapas, agrupando prefixos de casos
concluídos por fluxo de controle e encaminhando um caso em andamento, em tempo de execução, ao
classificador do grupo com que casa. Como a sequência é codificada é determinante de primeira ordem
do desempenho desses preditores \cite{teinemaa2019outcome}, e são justamente as técnicas voltadas a
modelos interpretáveis as que abrem mão de acurácia \cite{tax2020interdisciplinary}. As diferenças
deste trabalho em relação a essa linha são estreitas: casos concluídos inteiros em vez de prefixos,
cinco famílias de modelo em vez de uma, uma etapa separada que converte cada modelo por grupo em um
caminho inspecionável, e os dois controles acima no lugar da suposição de que o arranjo se paga. A
predição é aqui instrumento de medida, não produto.

Resumir um grupo por uma única execução tem maquinário formal por trás: o alinhamento múltiplo de
sequências expõe o núcleo conservado e os desvios \cite{bose2012diagnostics}, e há quem faça dos
traços do modelo os centroides dos grupos \cite{chatain2017alignment}, depois generalizados para um
submodelo \cite{boltenhagen2021modelbased}. Um caminho único é deslocamento deliberado para o canto
da simplicidade, comprado ao custo de generalização \cite{buijs2014quality}, e o grafo de precedência
direta aqui usado carrega distorções documentadas \cite{vanderaalst2019dfg}. O que difere é o que o método pode supor: centroides por
alinhamento exigem um modelo de processo como entrada, e a ausência de qualquer modelo documentado
do procedimento tal como ele é constitui a premissa deste estudo. O representante aqui adotado
deriva de um modelo sequencial treinado sobre os traços do próprio grupo, é comparável entre dois
alfabetos de traço e cinco famílias de modelo, e é pontuado contra o medoide --- o traço real mais
central do grupo ---, linha de base usual quando se pede um representante de grupo. O consenso por
alinhamento permanece o representante alternativo não computado.

A predição da próxima atividade é matéria madura, das redes recorrentes \cite{tax2017lstm} ao
tratamento de traços como linguagem \cite{evermann2017predicting} e a codificadores baseados em
atenção \cite{bukhsh2021processtransformer}. Os protocolos importam mais que as arquiteturas: sob
protocolo comum, a superioridade reportada frequentemente não se reproduz e nenhuma arquitetura
domina em todos os registros de atividades \cite{ramamaneiro2023dlppm}, e as partições em uso corrente vazam informação
entre treino e teste de modos que inflam o desempenho relatado \cite{weytjens2022leakage}. Um número
isolado de acurácia é afirmação sobre um registro de atividades e um protocolo, não sobre um método. Some-se a isso que
uma rede com boa acurácia de próxima atividade pode falhar em capturar a estrutura do processo
subjacente \cite{peeperkorn2023rnnstructure}: aderência e predição podem divergir, e uma não
substitui a outra sem argumento.

\subsection{Índices internos de validade de agrupamento}
\label{sec:rev-indices}

Os índices calculados neste trabalho foram introduzidos como instrumentos \emph{relativos}, isto é,
para ordenar partições candidatas produzidas em condições comparáveis, e não como medidas absolutas
de qualidade \cite{rousseeuw1987silhouettes,daviesbouldin1979measure,calinski1974dendrite,
moulavi2014dbcv}. Lidos como objetivos, comportam-se mal: partições que os maximizam diretamente
frequentemente discordam de rotulações de referência produzidas por especialistas
\cite{gagolewski2021invalid}; sobre 2,7 milhões de partições, \citeonline{yerbury2025rvi} os
concluem empiricamente inadequados para escolher entre abordagens de agrupamento, e recomendam
critérios orientados a resultado. Uma assimetria adicional incide sobre a comparação entre
algoritmos: nenhum índice consagrado avalia se as atribuições de ruído estão \emph{corretas}, de
modo que um agrupamento baseado em densidade e um particional não são pontuados sobre a mesma
população \cite{beer2026disco}. O gênero é familiar da modelagem de tópicos \cite{chang2009teahaves}
e do argumento de que o agrupamento não admite procedimento de qualidade independente de aplicação
\cite{vonluxburg2012clustering}. O contraexemplo é o de \citeonline{chicco2025dbcv}, que consideram
o DBCV mais informativo que três índices concorrentes baseados em densidade; o achado não é
contestado, mas compara índices quanto à recuperação da geometria dos grupos, com concordância de
60\% frente a um índice externo, e informatividade sobre geometria é grandeza distinta de poder
preditivo sobre utilidade a jusante.

A contribuição deste trabalho a esse debate é uma instância, e não uma teoria: um mesmo corpus
administrativo particionado de vários modos, com os índices e uma medida a jusante lado a lado.

\subsection{Privacidade em registros de atividades e mineração de processos no setor público}
\label{sec:rev-privacidade}

Registros de atividades são dados pessoais por direito próprio, porque um traço é impressão digital
comportamental que sobrevive à remoção de nomes \cite{mannhardt2019privacy,elkoumy2022privacy}. A
resposta da área tem sido a higienização desenhada em torno do que a análise a jusante precisa, com
o custo de utilidade medido: $k$-anonimato sobre variantes de traço
\cite{fahrenkrogpetersen2019pretsa}, extensão aos atributos contextuais dos eventos
\cite{fahrenkrogpetersen2020pripel}, uma formulação de otimização da fronteira entre privacidade e
utilidade \cite{fahrenkrogpetersen2023optimal}, o reagrupamento de uma coorte administrativa
anonimizada com três dos índices aqui utilizados \cite{ferrao2022utility} e a desidentificação de
texto para processamento de linguagem natural supervisionado \cite{lothritz2023evaluating}. Até onde
se pôde apurar, o efeito da desidentificação de texto livre sobre a qualidade do agrupamento não
supervisionado de traços e sobre a predição da próxima atividade não foi reportado, e é essa a
pergunta estreita a que este trabalho responde.

A mineração de processos já foi aplicada à administração pública e a tribunais, e à auditoria antes
de ambas \cite{vanderaalst2010auditing,jans2014fieldstudy}. \citeonline{unger2021jurimetrics}
constroem um registro de atividades a partir de ações de direito empresarial em tribunal estadual brasileiro e
constatam que os feitos seguem fluxo \emph{ad hoc}, e não desenhado;
\citeonline{nai2026procurement} reconstroem fluxos de contratação pública a partir de registros
europeus de licitação. A produção brasileira argumentou pela mineração de processos como instrumento
de transparência \cite{araujo2023mineracao} e analisou processos administrativos de um tribunal
regional eleitoral extraindo o registro de atividades do próprio SEI \cite{limaesilva2026analise} --- o mesmo sistema
de que provém o corpus deste trabalho ---, ao lado de estudos de caso em governo eletrônico e de
adoção em outras jurisdições \cite{delgado2022discovery,racis2025donot,trottier2025using}. Todos
analisam um procedimento cuja identidade é conhecida de antemão. A unidade de análise aqui é o
próprio conjunto de tipos de procedimento, porque a taxonomia declarada da administração não o
registra de forma confiável, e o catálogo daí resultante é o objeto que em seguida se submete a
avaliação.
